% page 94 du fac-simile (image p-093 du scan)
% bloc 160.0 x 233.3 mm ; marge gauche 8.0 mm
\pgc{-12.700mm}{0.000mm}{
\l{\cel{9}86}
\l{}
\l{\sou{chamaro}. Speco di redingoto kun kordoni transversa an-avane,}
\l{\cel{3}vesto Polona de epoko revolucionala di la XIX-esma yarcento}
\l{}
\l{}
\l{\sou{chambelano}. Oficiero komisita por okupar su pri la chambro di}
\l{\cel{3}rejo, di princo, qua +weris klefo kom insigno. - EFIS.}
\l{}
\l{\sou{chambro}. I. Parto di lojeyo. - II. Salono ube onu su asemblas}
\l{\cel{3}por deliberar. - III. En la guverno-sistemi reprezentala :}
\l{\cel{3}singla de la du asemblitari qui posedas la yuro legifar}
\l{\cel{3}same kam la chefo di la exekutigo-autoritato. - IV. (\sou{tek}-}
\l{\cel{3}\sou{nologio}).Nomo di kavaji diversa qui plu o min similesas}
\l{\cel{3}chambro. - EFI.}
\l{}
\l{\sou{champanio}. Vino de Champania (Francia). - DEFIRS.}
\l{}
\l{\sou{champiniono}. (\sou{bot.}) Genero de fungi di la familio \textquotedbl{}agaricinei\textquotedbl{},}
\l{\cel{3}karakterizata da la prezenteso di ringo ye la supro di la}
\l{\cel{3}pedo, lami libera rozea, pose nigra, t.e. departanta de}
\l{\cel{3}la bordo di la chapelo e qui ne arivas til la pedo,e spori}
\l{\cel{3}bruna purpurea. - F.}
\l{}
\l{\sou{championo}. Singla de laadversi qui kombatis en lico.(\sou{Metaf.})}
\l{\cel{3}La persono qua konsakras su a la defenso di ideo, di idealo.}
\l{\cel{3}- DEF.}
\l{}
\l{\sou{chanco}. Probableso favoroza o desfavoroza. - DEFI.}
\l{}
\l{\sou{chanfreno}. Parto avana di la kapo di kavalo, de la fronto til}
\l{\cel{3}la naztrui. - EF.}
\l{}
\l{\sou{chanjar}.\cel{2}(\sou{trans. e netrans.}) I. Substitucar ula kozo ad altra}
\l{\cel{3}en un loko o funciono. - II. Prenar kozo vice altra, o pos}
\l{\cel{3}altra. II. Divenar altra. - EFL.}
\l{}
\l{\sou{chankro}. (\sou{patol}.) Ulcero qua tendencas rodar la parti qui cir-}
\l{\cel{3}kondas lu. - DEFR.}
\l{}
\l{\sou{chanto}. Edro laterala, streta, di korpo larja e plata. - DF.}
\l{}
\l{\sou{chapo}. I. Mantelo ceremoniala quan +weras la episkopi, la}
\l{\cel{3}kantori, la sacerdoti por ula oficii solena. - II. (\sou{tekn.})}
\l{\cel{3}Peco qua parkovras ulo. - DFS.}
\l{}
\l{\sou{chapelo}. I. Kapovesto -virala) qua konsistas ordinare ye}
\l{\cel{3}korpo rigida o mi-rigida por recevar la kapo, kun rebordo}
\l{\cel{3}cirklo-kurva plu o min larja. - II. Kapovesto (mulierala)}
\l{\cel{3}qua konsistas generale ye korpo por recevar la kapo, o ka-}
\l{\cel{3}loto, kun rebordo. - FI.}
\l{}
\l{\sou{chapitro}. Singla de la parti ye qua konsistas la dividuri di}
\l{\cel{3}libro, kodexo, kontrato, budjeto, e c. - DEFIRS.}
}
